% ==============================================================================
%  4.2  Diagnóstico para el predictor X
% ==============================================================================
\section{Diagnóstico para el predictor}
\label{sec:t4-predictor}

El modelo no hace ninguna suposición específica sobre $X$, cuyos valores se consideran
conocidos (\cref{def:rls}). Por tanto, evaluar la idoneidad de $X$ como predictor no forma
parte de la validación del modelo propiamente dicha. Sin embargo, es necesario disponer de
información sobre $X$ para interpretar lo que ocurre con la respuesta. En concreto, interesa:
\begin{itemize}
  \item comprobar si hay \textbf{valores atípicos de $X$} que puedan influir en la recta de
    regresión ajustada (\cref{sec:t4-outliers-x});
  \item conocer el \textbf{rango y la dispersión} de los valores de $X$, para determinar el
    rango de validez de las predicciones realizadas con el modelo ajustado (el \emph{scope} del
    modelo, \cref{sec:t2-analisis}).
\end{itemize}
Para ello resulta útil un análisis descriptivo de $X$, con representaciones como diagramas de
cajas o histogramas.
